% Part: propositional-logic
% Chapter: syntax-and-semantics
% Section: formulas

\documentclass[../../../include/open-logic-section]{subfiles}

\begin{document}

\olfileid{pl}{syn}{fml}

\olsection{વિધાનાત્મક \usetoken{P}{formula}}

વિધાનાત્મક તર્કશાસ્ત્રનાં !!^{formula}s,
\emph{!!{propositional
    variable}s}માંથી\iftag{prvFalse}{\iftag{prvTrue}{,}{ અને}
  વિધાનાત્મક અચળ~$\lfalse$}{}\iftag{prvTrue}{\iftag{prvFalse}{
    અને}{} વિધાનાત્મક અચળ~$\ltrue$}{} \emph{તાર્કિક સંયોજકો} વડે
રચાય છે.

\begin{enumerate}
\item !!{propositional variable}s $\Obj p_0$, $\Obj p_1$, \dots નો
  !!^a{denumerable} ગણ~$\PVar$.
\tagitem{prvFalse}{!!{falsity} માટેનો વિધાનાત્મક અચળ~$\lfalse$.}{}
\tagitem{prvTrue}{!!{truth} માટેનો વિધાનાત્મક અચળ~$\ltrue$.}{}
\item તાર્કિક સંયોજકો:
  \startycommalist
  \iftag{prvNot}{\ycomma $\lnot$ (નિષેધ)}{}%
  \iftag{prvAnd}{\ycomma $\land$ (સંયોજન)}{}%
  \iftag{prvOr}{\ycomma $\lor$ (વિયોજન)}{}%
  \iftag{prvIf}{\ycomma $\lif$ (!!{conditional})}{}%
  \iftag{prvIff}{\ycomma $\liff$ (!!{biconditional})}{}%
\item વિરામચિહ્નો: (, ) અને અલ્પવિરામ.
\end{enumerate}

વિધાનાત્મક તર્કશાસ્ત્રની આ ભાષાને આપણે $\Lang L_0$થી દર્શાવીએ છીએ.

\iftag{defNot,defOr,defAnd,defIf,defIff,defTrue,defFalse,defEx,defAll}{%
ઉપર રજૂ કરેલા મૂળભૂત સંયોજકો ઉપરાંત, આપણે નીચેના \emph{વ્યાખ્યાયિત}
સંકેતો પણ વાપરીએ છીએ:
\startycommalist
  \iftag{defNot}{\ycomma $\lnot$ (નિષેધ)}{}%
  \iftag{defAnd}{\ycomma $\land$ (સંયોજન)}{}%
  \iftag{defOr}{\ycomma $\lor$ (વિયોજન)}{}%
  \iftag{defIf}{\ycomma $\lif$ (!!{conditional})}{}%
  \iftag{defIff}{\ycomma $\liff$ (!!{biconditional})}{}%
  \iftag{defFalse}{\ycomma $\lfalse$ (!!{falsity})}{}%
  \iftag{defTrue}{\ycomma $\ltrue$ (!!{truth})}}{}.

\begin{tagblock}{defNot,defOr,defAnd,defIf,defIff,defTrue,defFalse,defEx,defAll}
\begin{explain}
વ્યાખ્યાયિત સંકેત સત્તાવાર રીતે ભાષાનો ભાગ નથી, પણ અનૌપચારિક
સંક્ષેપ તરીકે રજૂ થાય છે: માત્ર મૂળભૂત સંકેતો વાપરીએ તો ઘણાં લાંબાં
થતાં સૂત્રોને તેના વડે ટૂંકાં લખી શકાય છે. આ દેખીતો લાભ છે. જોકે વધુ
મોટો લાભ એ છે કે સાબિતીઓ ટૂંકી થાય છે. કોઈ સંકેત મૂળભૂત હોય તો
સાબિતીઓમાં તેને અલગથી ધ્યાનમાં લેવો પડે. તેથી મૂળભૂત સંકેતો જેટલા
વધુ, આપણી સાબિતીઓ એટલી લાંબી.
\end{explain}
\end{tagblock}

% Alternate symbols

\begin{intro}
ઉપર આપણે વાપર્યાં તેનાથી જુદી પરિભાષા અને સંકેતોથી તમે પરિચિત હોઈ
શકો. તર્કશાસ્ત્રનાં પુસ્તકો (અને શિક્ષકો) ``નિષેધ'' માટે સામાન્ય રીતે
$\sim$, $\neg$ અને~!\; તથા ``સંયોજન'' માટે $\wedge$, $\cdot$ અને
$\&$માંથી કોઈ સંકેત વાપરે છે. ``શરતી વિધાન'' અથવા ``પ્રેરણ'' માટે
સામાન્ય સંકેતો $\rightarrow$, $\Rightarrow$ અને $\supset$ છે.
\iftag{prvIff,defIff}{``દ્વિશરતી વિધાન'', ``દ્વિમુખી પ્રેરણ'' અથવા
  ``(સત્યમૂલ્યલક્ષી) સમતુલ્યતા'' માટે $\leftrightarrow$,
  $\Leftrightarrow$ અને $\equiv$ વપરાય છે.}{}
\iftag{prvFalse,defFalse}{$\lfalse$ સંકેતને જુદી જગ્યાએ ``અસત્યતા'',
  ``ફાલ્સમ'', ``અસંગતિ'' અથવા ``તળિયું'' કહે છે.}{}
\iftag{prvTrue,defTrue}{$\ltrue$ સંકેતને જુદી જગ્યાએ ``સત્ય'',
  ``વેરમ'' અથવા ``ટોચ'' કહે છે.}{}
\end{intro}

\begin{defn}[સૂત્ર]
\ollabel{defn:formulas}
વિધાનાત્મક તર્કશાસ્ત્રનાં \emph{!!{formula}s}નો ગણ~$\Frm[L_0]$
નીચે પ્રમાણે અનુમાનાત્મક રીતે વ્યાખ્યાયિત થાય છે:
\begin{enumerate}
\tagitem{prvFalse}{$\lfalse$ એ આણ્વિક !!{formula} છે.}{}

\tagitem{prvTrue}{$\ltrue$ એ આણ્વિક !!{formula} છે.}{}

\item દરેક !!{propositional variable}~$\Obj p_i$ આણ્વિક
  !!{formula} છે.

\tagitem{prvNot}{જો $!A$ !!a{formula} હોય, તો $\lnot !A$ પણ
  !!a{formula} છે.}{}

\tagitem{prvAnd}{જો $!A$ અને $!B$ !!{formula}s હોય, તો $(!A \land
  !B)$ પણ !!a{formula} છે.}{}

\tagitem{prvOr}{જો $!A$ અને $!B$ !!{formula}s હોય, તો $(!A \lor !B)$
  પણ !!a{formula} છે.}{}

\tagitem{prvIf}{જો $!A$ અને $!B$ !!{formula}s હોય, તો $(!A \lif !B)$
  પણ !!a{formula} છે.}{}

\tagitem{prvIff}{જો $!A$ અને $!B$ !!{formula}s હોય, તો $(!A \liff !B)$
  પણ !!a{formula} છે.}{}

\tagitem{limitClause}{બીજું કશું !!a{formula} નથી.}{}
\end{enumerate}
\end{defn}

\begin{explain}
!!{formula}sની વ્યાખ્યા \emph{અનુમાનાત્મક વ્યાખ્યા} છે. મૂળભૂત રીતે,
આપણે !!{formula}sનો ગણ અનંત સંખ્યામાં તબક્કાઓમાં રચીએ છીએ. પહેલા
તબક્કામાં બધાં આણ્વિક સૂત્રોને સૂત્ર જાહેર કરીએ છીએ; આ વ્યાખ્યાના
પહેલા થોડા કિસ્સાઓ, એટલે કે
\iftag{prvTrue}{$\ltrue$, }{}%
\iftag{prvFalse}{$\lfalse$, }{}%
$\Obj p_i$ના કિસ્સાઓને અનુરૂપ છે. તેથી ``આણ્વિક !!{formula}''નો અર્થ
આ સ્વરૂપનું કોઈ પણ !!{formula} થાય છે.

વ્યાખ્યાના બીજા કિસ્સાઓ અગાઉથી રચેલાં !!{formula}sમાંથી નવાં
!!{formula}s રચવાના નિયમો આપે છે. બીજા તબક્કામાં આણ્વિક
!!{formula}sમાંથી !!{formula}s રચવા આપણે એ નિયમો વાપરી શકીએ છીએ.
ત્રીજા તબક્કામાં આણ્વિક સૂત્રો અને બીજા તબક્કામાં મળેલાં સૂત્રોમાંથી
નવાં સૂત્રો રચીએ છીએ, અને આમ આગળ ચાલે છે. આવા કોઈ તબક્કે છેવટે
રચાતી દરેક વસ્તુ, અને માત્ર એવી વસ્તુ, !!{formula} છે.
\end{explain}

દ્વિસ્થાની સંયોજક~$\ast$ વડે $!B$, $!C$માંથી રચેલું સૂત્ર
$(!B \ast !C)$ લખતી વખતે આપણે ઘણી વાર સૌથી બહારની કૌંસજોડી છોડી
દઈને માત્ર~$!B \ast !C$ લખીશું.

\begin{tagblock}{defNot,defOr,defAnd,defIf,defIff,defTrue,defFalse,defEx,defAll}
\begin{defn}
વ્યાખ્યાયિત કારકો વાપરીને રચેલાં સૂત્રો નીચે પ્રમાણે સમજવાનાં છે:

\begin{tagenumerate}{defTrue,defFalse,defNot,defOr,defAnd,defIf,defIff,defEx,defAll}

\tagitem{defTrue}{$\ltrue$ એ
  \iftag{prvFalse}{$\lnot\lfalse$}{$(!A \lor \lnot !A)$, જ્યાં
    $!A$ કોઈ નિશ્ચિત આણ્વિક !!{formula} છે}, તેનું સંક્ષિપ્ત રૂપ છે.}{}

\tagitem{defFalse}{$\lfalse$ એ
  \iftag{prvTrue}{$\lnot\ltrue$}{$(!A \land \lnot !A)$, જ્યાં
    $!A$ કોઈ નિશ્ચિત આણ્વિક !!{formula} છે}, તેનું સંક્ષિપ્ત રૂપ છે.}{}

\tagitem{defNot}{$\lnot !A$ એ $!A \lif \lfalse$નું સંક્ષિપ્ત રૂપ છે.}{}

\tagitem{defOr}{$!A \lor !B$ એ
  \iftag{prvAnd}{$\lnot(\lnot !A \land \lnot !B)$}{$\lnot !A \lif
    !B$}નું સંક્ષિપ્ત રૂપ છે.}{}

\tagitem{defAnd}{$!A \land !B$ એ
  \iftag{prvOr}{$\lnot(\lnot !A \lor \lnot !B)$}{$\lnot (!A \lif
    \lnot !B)$}નું સંક્ષિપ્ત રૂપ છે.}{}

\tagitem{defIf}{$!A \lif !B$ એ
  \iftag{prvOr}{$\lnot !A \lor !B$}{$\lnot (!A \land \lnot !B)$}નું
  સંક્ષિપ્ત રૂપ છે.}{}

\tagitem{defIff}{$!A \liff !B$ એ $(!A \lif !B) \land (!B
  \lif !A)$નું સંક્ષિપ્ત રૂપ છે.}{}
\end{tagenumerate}
\end{defn}
\end{tagblock}
\sourcecorrection{OLPL-001}{મૂળની \texttt{formulas.tex}માં વ્યાખ્યાયિત
પ્રેરણના વિયોજનવાળા વિકલ્પને $\lnot !A\lor !B)$ લખતાં એક વધારાનું
અંતકૌંસ આવે છે. અહીં અર્થાત્ બનતું સુઘડ સૂત્ર $\lnot !A\lor !B$
લખ્યું છે.}

\begin{defn}[વાક્યરચનાત્મક અભિન્નતા]
$\ident$ સંકેત, સંકેતોની શ્રેણીઓ વચ્ચેની વાક્યરચનાત્મક અભિન્નતા
દર્શાવે છે; એટલે કે, $!A \ident !B$ ત્યારે જ જ્યારે $!A$ અને $!B$
સમાન લંબાઈ ધરાવતી સંકેતશ્રેણીઓ હોય અને દરેક સ્થાને બંનેમાં એક જ
સંકેત હોય.
\end{defn}

$\ident$ સંકેતની બંને બાજુએ જોડાણથી મળેલી શ્રેણીઓ પણ આવી શકે છે.
દાખલા તરીકે, $!A \ident (!B \lor !C)$નો અર્થ આ છે: $!A$ની
સંકેતશ્રેણી એ આરંભકૌંસ, શ્રેણી $!B$, સંકેત $\lor$, શ્રેણી $!C$ અને
અંતકૌંસને આ જ ક્રમમાં જોડવાથી મળતી શ્રેણી છે. જો એવું હોય, તો આપણે
જાણીએ છીએ કે $!A$નો પહેલો સંકેત આરંભકૌંસ છે, $!A$માં $!B$ ઉપશ્રેણી
તરીકે આવે છે (બીજા સંકેતથી શરૂ થઈને), એ ઉપશ્રેણી પછી $\lor$ આવે છે,
વગેરે.

\end{document}
